\section{Introduction}
Hard disk drives (HDDs) are still the main back bones of modern storage system as HDDs are affordable for storing huge amount of data \cite{akhtar2019avic}. But compared to the rocky-rising request on cloud application, HDDs storage system is too slow for servicing highly demanded system \cite{armenatzoglou2022redshift}. A single disk can sustain perhaps 100 requests per second \cite{beckmann2018lhd}. This will dramatically deter the system's performance in peak accessing time. Considering that not all data are equally needed by users, cache is a cheap way to resolve this challenge. Especially when the flash memory is available, add a flash drive cache before HDDs storage is a good strategy to improve the storage system's performance \cite{berger2017adaptsize}. By storing most commonly accessed data in cache, the system will be able response to data request far more faster and it prevents most of the load from happening in rush hours on our hard drives. While the existing flash memory has a fatal weakness that it can only be write for a fixed times. When it's been overwritten for a certain times, it has to be replaced for failure. As it’s expensive to replace flash, balancing between caching and saving cost is a key success aspect to pursure for engineering.

To achieve a good cache performance, typically, cache systems uses admission policies to load data into it \cite{chakraborttii2021flash}. The admission policies determine what data and when to put in the flash cache. CoinFlip\cite{belady1966replacement} and RejectX\cite{berger2018lightweight} are commonly used cache adminission policies. They are effective for good hit-rate, but they are not designed to minimize the writes to flash cache. Most of them attempt to achieve a higher hit rate or minimize the flash used. This means they are optimizing the cache cost for the system. These techniques do not address the worst-case delays from the disk drives and the cache cost balancing.

Many researches are attempting to use AI strategies to optimize the cahce system, balancing the hit-rate and the cache replacement cost. Baleen \cite{wong2024baleen} is one of the latest work to address this. Wong, et, al used machine learning in their cache strategy to determine what data is worth caching. They grouped the disk data requests into groups which they referred as episodes, and using baleen to train their model to predict which ones shave the most time. Then in their work, they used train models based on past data from these episodes, to implement admission policies and data prefetch policies. It then employs the models in simulation to lower peak load on hard drives. They quantified the system performance by Peak Disk-head Time (Peak DT). In their work\cite{wong2024baleen}, the authors minimize the peak DT, by prefetch policies to load data into flash cache. According to their test result, the cache system has better performance than other methods because it takes into account the data hit-rate and the falsh writting. 

Convential flysh cache additi plicnies Such as ondin Coinrip and RejectX relyonrandomization, or simple frsguMc yheuritirs io res- aboutthama writes on the Fash-storedsWrite by-pass synd ducts \cite{berger2014ttlnetworks}. For instance, RejectX defers blocks in such a way that the block is not admitted to program until being read several times; this helps avoid wasteful writes and extend flash storage life \cite{berger2018lightweight}. These strategies are simple and can be easily applied, but they depend on strict thresholds and coarse access patterns that are not able to cope with workload variety. Importantly, they do so primarily by optimizing flash wear or hit rates and ignore system-level performance characteristics, such as maximum backend load or disk head-time.

To overcome these shortcomings, ML has been investigated as a solution to enhance cache-related decisions by researchers. The Flashield system~\cite{eisenman2019flashield} proposed a DRAM-buffered flash cache with an ML classifier to predict reuse using only the DRAM access information. Likewise, CacheLib-ML \cite{berg2020cachelib} that runs in Meta used gradient-boosted decision trees with frequency and metadata features to filter flash cache admission. But both are designed for IO hit rate or bandwidth efficiency—not for disk-head time (DT), which is the bottleneck of high-throughput HDD- backed storage systems. They also did not include any long-term access correlation or lifetime costbenefit tradeoff.

The gap is covered by Baleen system \cite{wong2024baleen}, which presents a new cache residency model via episodes to better model access clusters across time. Baleen trains policy (supervised ML) to mimic an oracle admission pol-icy (OPT) which minimizes DT subject to flash write rate limits. In contrast with previous approaches, Baleen explicitly optimizes for Peak DT, which directly corresponds to the backend provisioning requirement. Furthermore, Baleen has been designed in modular form with decoupled admission, prefetching, and eviction modules making it feasible for a collective learning approach all the while maintaining explainability and system performance. In experimental testing on production traces, Baleen decreased peak backend load by 16\% and TCO for both immediately-served and backfilled requests by 17\% compared to the best baseline. 

Baleen already gives a strong base for cache research, but while working on it we still noticed some parts that feel not fully done yet. A lot of older works pay more attention to admission and prefetching, but not so much on eviction~\cite{wong2024baleen}. Things like parameter sensitivity or tuning were mentioned before, but not really tested deeply. Because of that, when workload changes or cache size becomes different, those systems sometimes behave weird or unstable. We think this kind of issue makes the system less flexible. So our main goal is to improve Baleen and try to make it perform more steady across different situations.

The idea is basically to let the system react faster when backend load changes. We mainly have two parts in our design: DT-SLRU (E1) and Episode-Deadline Eviction (E2). DT-SLRU changes its promotion threshold based on disk service time~\cite{beckmann2018lhd}, trying to find some balance between short-term reaction and long-term reuse. Episode-Deadline Eviction tries to guess how long each block might stay useful and gives it a soft “deadline” to remove it. It also keeps a small protected area for blocks with high DT-per-byte, since those are usually worth keeping longer. From what we’ve seen in early testing, the setup seems smoother and more consistent when workload becomes unstable.

Overall, our design shows that mixing ML-based guidance with temporal eviction can help the cache system stay stable and more adaptive to different workloads. It’s still not perfect, but already looks more reliable, especially when things get unpredictable or heavy.

Our contributions are summarized as follows:
\begin{itemize}[topsep=0pt, parsep=0pt, itemsep=0pt, leftmargin=*]
    \item We reproduce and extend the Baleen (FAST~'24) framework, focused on the underdeveloped eviction components in the ML-guided caching system.
    \item We implemented two adaptive eviction schemes --- DT-SLRU and Episode-Deadline Eviction (EDE).
    \item We performed ablation studies by adjusting the promotion threshold ($\tau_{DT}$) and protected-cap ratio, and examined how these parameters will affect Peak DT, median DT, hit rate, and flash write traffic.
    \item Our experiments shows that proper tuning of these parameters could lower Peak DT by around 45\% comparing to the baseline.
\end{itemize}
